% =====================================================================
%  MLE-AI-201 Machine Learning -- Final Project report (IEEE-style, self-contained)
%
%  EVERY number, table and figure below is produced by
%        python -m src.run_all
%  which writes  report/generated/macros.tex, report/generated/tab_*.tex
%  and           report/figures/*.pdf.   Build afterwards with
%        cd report && pdflatex report.tex && pdflatex report.tex
%  (or `make report`).  Do not type results by hand.
% =====================================================================
\documentclass[10pt,twocolumn,a4paper]{article}

\usepackage[a4paper,top=0.75in,bottom=1.0in,left=0.68in,right=0.68in,columnsep=0.25in]{geometry}
\usepackage{mathptmx}
\renewcommand{\ttdefault}{lmtt}   % vector typewriter font in T1 (Latin Modern)
\usepackage[T1]{fontenc}
\usepackage[utf8]{inputenc}
\usepackage{graphicx}
\usepackage{amsmath,amssymb}
\usepackage{booktabs}
\usepackage{array}
\usepackage{titlesec}
\usepackage{caption}
\usepackage{cite}
\usepackage[hidelinks]{hyperref}
\usepackage{url}
\usepackage{xcolor}

\renewcommand{\thesection}{\Roman{section}}
\renewcommand{\thesubsection}{\Alph{subsection}}
\titleformat{\section}[block]{\centering\normalsize\scshape}{\thesection.}{0.5em}{}
\titleformat{\subsection}[block]{\normalsize\itshape}{\thesubsection.}{0.5em}{}
\titlespacing*{\section}{0pt}{1.4ex plus 1ex minus .2ex}{1ex}
\titlespacing*{\subsection}{0pt}{1.0ex plus 1ex minus .2ex}{0.6ex}
\newcommand{\abstractkw}[1]{\noindent\textbf{\textit{#1}}}
\captionsetup{font=small,labelsep=period}
\setlength{\parindent}{1em}
\graphicspath{{figures/}}
% let LaTeX put more/larger floats on a page (many full-width figures)
\renewcommand{\topfraction}{0.92}\renewcommand{\bottomfraction}{0.8}
\renewcommand{\textfraction}{0.06}\renewcommand{\floatpagefraction}{0.75}
\renewcommand{\dbltopfraction}{0.92}\renewcommand{\dblfloatpagefraction}{0.75}
\setcounter{topnumber}{3}\setcounter{dbltopnumber}{3}\setcounter{totalnumber}{6}

% ---- generated results (run `python -m src.run_all` first) -------------
\IfFileExists{generated/macros.tex}{\input{generated/macros.tex}}{%
  \errmessage{report/generated/macros.tex is missing: run `python -m src.run_all` from the repo root first}}
\newcommand{\gen}[1]{\input{generated/#1.tex}}
% shrink a generated table to the line width only if it is wider than that
\newcommand{\fit}[1]{\resizebox{\ifdim\width>\linewidth\linewidth\else\width\fi}{!}{#1}}
\newcommand{\repo}{\href{https://github.com/eyyubovsamur240-afk/ML-final-project}{\texttt{github.com/eyyubovsamur240-afk/ML-final-project}}}
\newcommand{\vect}[1]{\mathbf{#1}}
\newcommand{\schwa}{\raisebox{0.13ex}{\rotatebox[origin=c]{180}{e}}}  % Azerbaijani schwa without extra packages

% ---- TEAM: fill in (alphabetical by surname) ----------------------------
\newcommand{\TeamAuthors}{Surname1, Name1\quad Surname2, Name2\quad Surname3, Name3}

% =====================================================================
\begin{document}

\twocolumn[
  \begin{@twocolumnfalse}
  \begin{center}
    {\LARGE\bfseries Decision Trees and Pegasos SVMs from Scratch:\\
     Price and Price-Tier Prediction on bina.az Listings\par}
    \vspace{1.2em}
    {\large \TeamAuthors\par}
    \vspace{0.5em}
    {\itshape MLE-AI-201 Machine Learning, AI Academy --- Cohort II\par}
    \vspace{1.4em}
  \end{center}
  \begin{quote}
    \small
    \abstractkw{Abstract---}%
    We implement a CART decision tree (Gini, entropy and MSE criteria) and a
    soft-margin support-vector machine trained with the Pegasos sub-gradient
    method entirely in NumPy, and apply them to \NClean{} cleaned sale
    listings scraped from bina.az, Azerbaijan's largest real-estate portal.
    One dataset is framed as two tasks: regression of $\log(\text{price})$
    (Task~A) and classification of a \emph{premium} vs.\ \emph{standard}
    price tier split at the training median (Task~B). Every hyperparameter
    is selected on a validation split; final models are refitted on
    train+validation and scored once on a held-out test set, with
    \CVFolds-fold cross-validation for uncertainty. Our tree reaches a test
    RMSE of \OurTreeRMSE{} on $\log(\text{price})$ ($R^2=\OurTreeRtwo$),
    against \SkTreeRMSE{} for scikit-learn's tree with matched settings; a
    node-by-node comparison leaves \RegGainMismatch{} unexplained
    differences, with the remaining ones due to \TreeGapCause.
    Our averaged Pegasos SVM reaches test F1 \OurSVMFOne{}
    (scikit-learn LinearSVC: \SkLinSVCFOne), and a random-Fourier-feature
    extension reaches \OurRFFFOne. Error analysis shows that misclassified
    tiers concentrate around the threshold, and that the regression tree
    \TopDecileWord{} the most expensive listings.
    Code and one-command reproduction: \repo{} (tag \texttt{v1.0-final}).

    \vspace{0.6em}
    \abstractkw{Index Terms---}decision tree, support-vector machine,
    Pegasos, random Fourier features, real-estate price prediction.
  \end{quote}
  \vspace{1.0em}
  \end{@twocolumnfalse}
]

% =====================================================================
\section{Introduction}
Asking prices on listing portals are noisy, heavy-tailed and driven by
location, size and condition. We use them to study two classic learners
from the inside: we write both models ourselves, feed them a real,
uncleaned scrape, and benchmark them honestly against scikit-learn.
Our contributions are:
\begin{itemize}\setlength\itemsep{0pt}
  \item a leakage-free pipeline for the bina.az dump: tolerant parsing of
        Azerbaijani, unit-laden fields, currency conversion, documented
        cleaning rules, and preprocessing fitted on training rows only;
  \item a vectorised CART tree whose split search evaluates every
        threshold of every feature with cumulative sums, verified to make the
        same split as scikit-learn at every node up to exact ties;
  \item a Pegasos SVM with mini-batches, projection, iterate averaging and
        a random-Fourier-feature (RFF) kernel extension, whose objective we
        compare with LIBLINEAR's exact optimum;
  \item a protocol that never touches the test split during selection, CV
        uncertainty for all models, error analysis, and bonus ensembles,
        ridge regression and clustering built on our own code.
\end{itemize}

\section{Background}
\emph{Decision trees.} CART~\cite{cart} grows a binary tree greedily:
each node picks the feature/threshold pair that maximises the impurity
decrease $\Delta I = I(\text{parent})-\sum_{c}\frac{N_c}{N}I(c)$, where
$I$ is the Gini index $1-\sum_k p_k^2$, the entropy $-\sum_k p_k\log_2 p_k$
(information gain~\cite{quinlan}), or, for regression, the node variance
$\frac{1}{N}\sum_i (y_i-\bar y)^2$. Growth stops on depth, node-size or
minimum-gain rules; leaves predict the majority class (with class
frequencies as probabilities) or the mean.

\emph{SVMs and Pegasos.} The soft-margin SVM~\cite{cortes} minimises the
regularised hinge loss. Pegasos~\cite{pegasos} solves this primal by
stochastic sub-gradient descent with step $\eta_t=1/(\lambda t)$ and
reaches $\varepsilon$ accuracy in $\tilde O(1/(\lambda\varepsilon))$
steps, independent of $n$. Random Fourier features~\cite{rahimi} map
$\vect{x}$ to $\phi(\vect{x})=\sqrt{2/D}\cos(W\vect{x}+\vect{c})$ with
$W\sim\mathcal N(0,2\gamma I)$, so that
$\phi(\vect{x})^{\top}\phi(\vect{z})\approx e^{-\gamma\lVert\vect{x}-\vect{z}\rVert^2}$
and a linear solver on $\phi$ approximates an RBF-kernel SVM.

% =====================================================================
\section{Data and Preprocessing}
\subsection{Dataset and cleaning}
We use the Kaggle \emph{Bina.az sale project} dump~\cite{data} (file
\texttt{\DataFile}, SHA-256 prefix \texttt{\DataSHA}), \NRaw{} rows. Column
names are Azerbaijani (\emph{Sah\schwa{}}, \emph{Otaq say\i},
\emph{M\schwa{}rt\schwa{}b\schwa{}}, \ldots), most numbers arrive as
strings with units (``85~m$^2$'', ``5 / 9'', ``1~250~000'') and prices are
quoted in several currencies (\CurrencyBreakdown). All cleaning lives in
\texttt{src/data\_prep.py}:
\begin{enumerate}\setlength\itemsep{0pt}
  \item names are transliterated and mapped to English
        (\emph{area, rooms, floor, land area, repair, mortgage, bill of sale,
        category, building type});
  \item tolerant parsers handle thousands separators, decimal commas,
        ``floor / total floors'', land in \emph{sot}/m$^2$/ha and
        \emph{var}/\emph{yoxdur} flags; prices are converted to AZN at fixed
        rates (USD at the official 1.70 peg);
  \item rows without a positive price or area are dropped, then exact
        duplicates and re-posts (same price, area, rooms, floor and
        coordinates);
  \item fixed plausibility rules remove typos and unit mix-ups
        (Table~\ref{tab:cleaning}); impossible feature values (floor above
        the building height, coordinates outside Azerbaijan) are set to
        missing and imputed later.
\end{enumerate}
The rules are constants chosen from domain knowledge, not statistics
estimated on the data, so applying them before the split cannot leak.

\begin{table}[t]
  \centering
  \caption{Cleaning steps and rows remaining.}
  \label{tab:cleaning}
  \fit{\gen{tab_cleaning}}
\end{table}

\emph{Leakage.} We drop \textbf{\LeakageCols} because they are
arithmetic functions of the target (price per m$^2$, price restated);
the cleaning code also drops any other column whose name contains
``price''. We also drop identifiers and free-form addresses
(\IdentifierCols) and platform promotion flags (\MetaCols), which describe
the seller's advertising budget rather than the property. Price per m$^2$
is used once, only as a typo filter, and never as a feature.

\subsection{Features, tasks and split}
From the cleaned rows we build \NNumericFeatures{} numeric/binary features
(area and $\log$ area, rooms, area per room, floor, total floors, floor
ratio, top/ground-floor flags, land area, latitude, longitude, distance
to the Baku centre, repair/mortgage/bill-of-sale flags and six keyword
flags from the description, e.g.\ ``metro'' or ``d\schwa{}niz'' (sea)), plus
one-hot encodings of category, building type, city and district
(\NFeatures{} columns in total). One-hot rather than ordinal encoding is
used because none of these has a natural order. Missing numeric values are
imputed with the \emph{training} median and flagged by an indicator, and
a category level gets its own column only if it has at least 20
\emph{training} rows; rarer or unseen levels fall into an ``other'' column.

Task~A predicts $\log(\text{price})$: the raw price is strongly
right-skewed (skewness \PriceSkew; \LogPriceSkew{} after the log,
Fig.~\ref{fig:eda}a--b), and an error in $\log$ space is a relative error.
Task~B labels a listing \emph{premium} if its price exceeds the
\emph{training} median, \TierThreshold{} AZN. The split is 70/15/15
(\NTrain/\NVal/\NTest), stratified on deciles of $\log(\text{price})$ with
a fixed seed, so all splits share the price distribution and the tier
balance (\PremiumTrain, \PremiumVal, \PremiumTest{} premium). Stratifying
only decides which rows go where; no statistic is computed from
validation or test rows.

\begin{table}[t]
  \centering
  \caption{Dataset summary after cleaning.}
  \label{tab:data}
  \fit{\gen{tab_dataset}}
\end{table}

\begin{figure*}[tp]
  \centering
  \includegraphics[width=\textwidth]{eda_price_distribution}\\[2pt]
  \includegraphics[width=0.49\textwidth]{eda_price_map}\hfill
  \includegraphics[width=0.49\textwidth]{eda_area_vs_price}
  \caption{Exploration on the training split. (a, b) Price is heavy-tailed;
  $\log(\text{price})$ is close to symmetric. (c) Listing prices across
  Baku: the centre and the seafront are the most expensive.
  (d) $\log$ area vs.\ $\log$ price (correlation \CorrAreaPrice).}
  \label{fig:eda}
\end{figure*}

\emph{Exploration} (training split only, Fig.~\ref{fig:eda} and the
appendix). The median price is \PriceMedian{} AZN (1st--99th percentile
\PricePOne--\PricePNinetyNine), the median price per m$^2$ is \PpmMedian{}
AZN, and $\log$ area correlates with $\log$ price at \CorrAreaPrice. The
median price per m$^2$ varies several-fold between districts
(Fig.~\ref{fig:loc}), so location carries a large part of the signal.
After cleaning the most frequent missing fields are \MissingSummary.

% =====================================================================
\section{Method}
\subsection{Decision tree (from scratch)}
\texttt{DecisionTree} (\texttt{src/decision\_tree.py}) builds \texttt{Node}
objects recursively and exposes a scikit-learn-style
\texttt{fit/predict/predict\_proba} API. \emph{Split search:} for a node
with $m$ rows we sort all candidate features at once (one
\texttt{argsort} of the $m\times d$ block) and evaluate \emph{every}
threshold with cumulative sums. For classification, running class counts
give the Gini or entropy of both children at every cut. For regression,
running $\sum y$ and $\sum y^2$ give the children's sums of squared
errors, computed after centring $y$ at the node mean for numerical
stability. Candidate thresholds are midpoints between consecutive
\emph{distinct} sorted values. A split is valid only if both children keep
at least \texttt{min\_samples\_leaf} rows, and the best gain wins (ties go
to the lowest feature index). One node therefore costs
$O(d\,m\log m)$, almost all of it inside NumPy.

\emph{Stopping and regularisation.} A node becomes a leaf if it is pure,
reaches \texttt{max\_depth}, has fewer than \texttt{min\_samples\_split}
rows, has no valid cut, or if the weighted gain
$\frac{N_t}{N}\Delta I$ is below \texttt{min\_impurity\_decrease}. This is
the same weighted definition scikit-learn uses, so matched settings mean
the same thing in both libraries. Leaves store class frequencies (giving
\texttt{predict\_proba}) or the mean target. Feature importance is the
total weighted gain per feature, normalised to one (mean decrease in
impurity). After fitting, the tree is flattened into arrays so that
prediction routes all rows through the tree together.

\emph{Depth curves from a single fit.} Greedy CART decides each node from
that node's rows only, so a tree grown to depth $D$ and \emph{truncated}
at depth $d$ is identical to a tree grown with \texttt{max\_depth}$=d$. A
unit test checks this equivalence, and we use it to draw depth curves at
the cost of one fit.

\subsection{Support-vector machine (from scratch)}
\texttt{PegasosSVM} (\texttt{src/svm.py}) minimises
\begin{equation}
  J(\vect{w})=\tfrac{\lambda}{2}\lVert\vect{w}\rVert^{2}
  +\tfrac{1}{n}\sum_{i=1}^{n}\max\!\big(0,\,1-y_i\,\vect{w}^{\top}\vect{x}_i\big),
  \quad y_i\in\{-1,+1\},
\end{equation}
where $\vect{x}_i$ is augmented with a constant 1 so that the bias is the
last weight. With a random mini-batch $A_t$ of size $k$ and the set
$A_t^{+}$ of its hinge-active points ($y_i\vect{w}^{\top}\vect{x}_i<1$),
one step is
\begin{equation}
  \vect{w}\leftarrow(1-\eta_t\lambda)\,\vect{w}+\frac{\eta_t}{k}\sum_{i\in A_t^{+}}y_i\vect{x}_i,
  \qquad \eta_t=\frac{1}{\lambda t},
\end{equation}
followed by the optional projection onto the ball
$\lVert\vect{w}\rVert\le1/\sqrt\lambda$, which contains the optimum.
$k=1$ is classic Pegasos and $k=n$ is full-batch sub-gradient descent;
the $\eta_0/\sqrt t$ and constant schedules are also implemented. Each
epoch reshuffles the rows with a seeded generator and records
$J(\vect{w})$ on the full training set. Because the last iterate is noisy
when $\lambda$ is small, the tuned solver returns the $t$-weighted iterate
average
$\bar{\vect w}_t=(1-\frac{2}{t+1})\bar{\vect w}_{t-1}+\frac{2}{t+1}\vect w_t$,
which keeps the optimal $O(1/(\lambda T))$ rate~\cite{lacoste}. Features
are standardised with training statistics, since the hinge loss is
scale-sensitive. We report $\lVert\vect w\rVert$, the margin width
$2/\lVert\vect w\rVert$, and the support vectors, defined as the training
points with $y_i f(\vect{x}_i)\le 1$.

\emph{Non-linear extension: random Fourier features.} We map the
standardised features with $D=\RFFD$ random features and run the
\emph{same} linear solver on $\phi(\vect x)$ (\texttt{RFFPegasosSVM}). We
chose RFF over kernelised Pegasos because training and prediction cost
$O(nD)$ instead of growing with the number of support vectors, which
matters for $\sim$\NDev{} training rows, and because it reuses the tested
linear code unchanged. The cost is a Monte-Carlo kernel approximation,
whose error we measure (Fig.~\ref{fig:svmall}, bottom-c).

% =====================================================================
\section{Experimental Setup}
\emph{Protocol.} (1) All hyperparameters are chosen on the validation
split with models trained on the training split; the choice rule is fixed
in code (best validation score, ties to the simpler model).
(2) The chosen configurations are refitted on train+validation
(\NDev{} rows) with imputation, encoding, scaling and the tier threshold
re-learned on those rows (threshold \TierThresholdDev{} AZN), and scored
once on the test split with 95\% percentile-bootstrap intervals (1{,}000
resamples). (3) For uncertainty, \CVFolds-fold stratified CV on
train+validation repeats the whole preprocessing inside every fold and
reports mean$\pm$std.

\emph{Metrics.} Task~A: RMSE, MAE and $R^2$ on $\log(\text{price})$, plus
MAE (AZN) and MAPE after back-transforming with $\exp$, which predicts the
conditional median. Task~B: accuracy, precision, recall and F1 of the
premium class, ROC-AUC and PR-AUC (average precision) from continuous
scores (leaf frequencies for trees, margins for SVMs), and confusion
matrices. All metrics are our own NumPy code and are unit-tested against
scikit-learn.

\emph{Baselines with matched settings.} Our tree vs.\
\texttt{DecisionTreeRegressor/Classifier} with the same criterion,
\texttt{max\_depth}, \texttt{min\_samples\_leaf} and
\texttt{min\_impurity\_decrease}. Our SVM vs.\ \texttt{SGDClassifier}
(hinge, $\alpha=\lambda$, ``optimal'' schedule~\cite{bottou}) and
\texttt{LinearSVC} ($C=1/(n\lambda)$, LIBLINEAR~\cite{liblinear}), both
behind a \texttt{StandardScaler} inside a \texttt{Pipeline}. Our RFF SVM
vs.\ the exact-kernel \texttt{SVC} with the same $\gamma$ and $C$. A
constant predictor (training median) anchors Task~A. All randomness is
seeded (\texttt{SEED=42}).

% =====================================================================
\section{Results}
\subsection{Tree: criteria, depth and the bias--variance trade-off}
Fig.~\ref{fig:depth} shows the classic picture. Training error keeps
falling with depth, while validation error bottoms out and then rises
again as the tree starts fitting noise. An unpruned regression tree
reaches a training RMSE of \TreeDeepTrainRMSE{} at depth \TreeDeepDepth{},
but its validation RMSE is \TreeDeepValRMSE; the best unpruned depth is
\TreeBestDepthUnpruned. Gini and entropy \CritSimilarity{} on
Task~B: their best validation F1 differ by \CritDiff{} (Gini \GiniValFOne,
entropy \EntropyValFOne). Larger leaves regularise a fully grown tree, and
\texttt{min\_impurity\_decrease} prunes it: at \MinDecLast{} the tree
shrinks from \MinDecNodesFirst{} to \MinDecNodesLast{} nodes
(Fig.~\ref{fig:tree-extra}). On validation we select
\textbf{\TreeClfCrit, depth \TreeClfDepth, min\_samples\_leaf
\TreeClfLeaf} for Task~B and \textbf{depth \TreeRegDepth,
min\_samples\_leaf \TreeRegLeaf, min\_impurity\_decrease \TreeRegMinDec}
for Task~A (validation RMSE \TreeRegValRMSE). In the learning curve
(appendix) the train--validation gap goes from \LCGapSmall{} to
\LCGapFull{} as data grows, and \LCVerdict.

\begin{figure*}[tp]
  \centering
  \includegraphics[width=\textwidth]{tree_depth_curves}
  \caption{Depth sweep with \texttt{min\_samples\_leaf}=1 (dotted: train,
  solid: validation). Validation error is U-shaped, the bias--variance
  trade-off. The grey line marks the depth chosen on validation.}
  \label{fig:depth}
\end{figure*}

\subsection{SVM: regularisation, margin and convergence}
Table~\ref{tab:svm} and Fig.~\ref{fig:svmall} (top) trace the effect of
$\lambda$. A larger $\lambda$ shrinks $\lVert\vect w\rVert$, which widens
the margin $2/\lVert\vect w\rVert$ and makes more training points fall
inside it as support vectors, trading training fit for a simpler
boundary. Validation selects $\lambda=\SVMLambda$ ($C=\SVMC$). At this
setting the final model has $\lVert\vect w\rVert=\SVMWNorm$, a margin of
\SVMMargin{}, and \SVMNSV{} support vectors (\SVMSVFrac{} of the training
rows).

Fig.~\ref{fig:svmall} (middle) confirms convergence: $J$ falls from \ObjFirst{} at
$\vect w=0$ to \ObjLast{} and changes by \ObjRelChange{} over the last five
epochs. Small $\lambda$ converges slowly, as the $1/(\lambda T)$ rate
predicts, which is why averaging matters. With the same epoch budget,
mini-batches of \SVMBatch{} \BatchVerdict{} $k=1$ at a fraction of the
wall time (\BatchTimes), while full-batch descent, which takes only one
step per epoch, \FullBatchVerdict. On the full training set,
our averaged Pegasos ends at $J=\ObjPegasos$, against LIBLINEAR's exact
optimum $J=\ObjLiblinear$ (gap \ObjGap), and the two agree on
\SVMLibAgree{} of test predictions.

\begin{table}[t]
  \centering
  \caption{Linear Pegasos SVM vs.\ $\lambda$ (train$\to$validation; the
  selected row is bold). $C=1/(n\lambda)$.}
  \label{tab:svm}
  \fit{\gen{tab_svm_sweep}}
\end{table}

\begin{figure*}[tp]
  \centering
  \includegraphics[width=\textwidth]{svm_lambda_sweep}\\[1pt]
  {\small\textbf{Top:} effect of $\lambda$ on (a) validation F1/AUC, (b) the margin width
  $2/\lVert\vect w\rVert$, (c) the share of support vectors.}\\[4pt]
  \includegraphics[width=\textwidth]{svm_convergence}\\[1pt]
  {\small\textbf{Middle:} training objective per epoch for (a) different $\lambda$,
  (b) batch size and iterate averaging at the selected $\lambda$, (c) step-size schedules.}\\[4pt]
  \includegraphics[width=\textwidth]{svm_rff_study}
  \caption{Our Pegasos SVM. \textbf{Bottom:} RFF extension: (a) bandwidth $\gamma\times\lambda$,
  (b) validation AUC vs.\ number of random features $D$, (c) mean absolute error of the kernel
  approximation.}
  \label{fig:svmall}
\end{figure*}



For the RFF extension, validation selects $\gamma=\RFFGamma$ and
$\lambda=\RFFLambda$ (Fig.~\ref{fig:svmall}, bottom). The kernel-approximation error
falls with a fitted log--log slope of \RFFErrSlope{} (theory: $-0.5$,
i.e.\ $1/\sqrt{D}$), and validation AUC \RFFDVerdict.



\subsection{Benchmark against scikit-learn}
Tables~\ref{tab:reg}--\ref{tab:headline} are the main results.
\emph{Task~A:} our tree's test RMSE is \OurTreeRMSE{} (CV
\CVOurTreeRMSE), scikit-learn's is \SkTreeRMSE{} (CV \CVSkTreeRMSE), and
the constant baseline's is \ConstRMSE. The back-transformed error is
\OurTreeMAEAZN{} AZN MAE (MAPE \OurTreeMAPE). The bonus ridge regression
(\RidgeRMSE) shows how much of the signal is close to linear in
$\log$ space: the single tree \TreeVsRidgeWord{} it. \emph{Task~B:} the
best of our models is \BestOurClf{} (F1 \BestOurClfFOne),
\BestOverallClause. The better of our two SVMs \SVMVsTreeWord{} our tree,
and the RFF kernel \RFFVsLinearWord{} the linear SVM (test F1
\RFFMinusLinear).
Averaged Pegasos is within \PegasosVsLiblinearFOneDiff{} F1 of LIBLINEAR.

\begin{table*}[tp]
  \centering
  \caption{Task A, regression of $\log(\text{price})$. CV:
  \CVFolds-fold mean$\pm$std on train+validation. Test: single held-out
  split with a 95\% bootstrap CI. MAE (AZN) and MAPE after
  back-transformation. Best test RMSE in bold.}
  \label{tab:reg}
  \fit{\gen{tab_results_reg}}
  \vspace{6pt}
  \caption{Task B, price tier (premium = positive class). Same protocol
  as Table~\ref{tab:reg}. Best test F1 in bold. $^\dagger$\SVCNote.}
  \label{tab:clf}
  \fit{\gen{tab_results_clf}}
\end{table*}



\begin{table}[t]
  \centering
  \caption{Headline results on both tasks (Task A: RMSE of
  $\log$ price, lower is better; Task B: F1, higher is better).}
  \label{tab:headline}
  \fit{\gen{tab_headline}}
\end{table}

\emph{Structure.} Table~\ref{tab:tree-compare} compares the fitted trees
node by node (\texttt{src/tree\_compare.py}). At each node we check
whether both libraries split the same training rows. If they do not, we
classify the difference: an exact tie (two splits with equal gain;
scikit-learn breaks ties by visiting features in random order), a
\emph{float32 divergence} (scikit-learn casts $X$ to \texttt{float32},
which here merges \FloatCollapsed{} distinct \texttt{float64} values; we
re-run our own split search on the \texttt{float32}-rounded rows of that
node and it returns exactly scikit-learn's split), or an unexplained
difference. We find \RegGainMismatch{} (Task~A) and \ClfGainMismatch{}
(Task~B) unexplained differences, \RegTies{} ties and \RegFloatDiv{} float32
divergences in Task~A, and test-prediction agreement of \RegAgreement{}
and \ClfAgreement. \RootSplitSentence

\emph{Speed.} Our trees fit in \OurTreeRegFit{}\,s vs.\
\SkTreeRegFit{}\,s for scikit-learn's Cython code. Our split search is
asymptotically the same $O(d\,m\log m)$ per node, but it pays Python
overhead per node and re-sorts at every node instead of reusing a
presorted order. Pegasos costs $O(\text{epochs}\cdot n d)$ and fits in
\OurSVMFit\,s. The exact-kernel SVC is $O(n^2)$--$O(n^3)$ (\SkSVCFit\,s),
the main practical argument for RFF.

\begin{table}[t]
  \centering
  \caption{Our tree vs.\ scikit-learn, node by node (final models; depth
  and leaves as ours / sklearn). Same: identical split. Ties: equal-gain
  alternatives. f32: explained by scikit-learn's float32 cast. Other:
  unexplained (0 for a correct implementation).}
  \label{tab:tree-compare}
  \fit{\gen{tab_tree_compare}}
\end{table}

\subsection{Error analysis}
\emph{What the models use.} Our tree's impurity-based importances
(Fig.~\ref{fig:imp}) are led by \TopFeatReg{} for Task~A and by
\TopFeatClf{} for Task~B. Size and location dominate, which matches the
EDA. In the linear SVM (Fig.~\ref{fig:weights}), whose weights are comparable because the
features are standardised, the largest positive weights are
\SVMPosWeights{} and the most negative are \SVMNegWeights.

\emph{Where the regression fails.}
(i) \emph{Luxury and bargain outliers.} The tree \TopDecileWord{} the top
price decile (mean $\log$ residual \TopDecileBias) and
\BottomDecileWord{} the bottom one (\BottomDecileBias,
Fig.~\ref{fig:err}, bottom left). \ShrinkageSentence{}
Table~\ref{tab:worst} lists the largest errors.
(ii) \emph{Sparse districts} (Fig.~\ref{fig:err}, top middle).
\SparseDistrictSentence{}
(iii) \emph{Categories.} The error is highest for \WorstCategory{}
(\WorstCategoryMAE) and lowest for \BestCategory{} (\BestCategoryMAE).
\ShareBigErrors{} of test listings are off by more than 50\% in price.

\emph{Where the classifier fails.} Tier errors concentrate at the
boundary. \NearThresholdErrShare{} of the tree's misclassifications (SVM:
\NearThresholdSVMErrShare) are listings within $\pm$10\% of the
threshold price, a band that holds only \NearThresholdRowShare{} of the
test rows. The error rate falls from \NearestThresholdErr{} closest to
the threshold to \FarThresholdErr{} far from it (Fig.~\ref{fig:err},
bottom right).
Many of these ``errors'' are therefore near-ties that no feature could
resolve: a listing just above the median is not meaningfully more
premium than one just below it.

\begin{figure*}[tp]
  \centering
  \includegraphics[width=\textwidth]{feature_importance}
  \caption{Top impurity-decrease (MDI) importances of our final trees.}
  \label{fig:imp}
\end{figure*}

\begin{figure*}[tp]
  \centering
  \includegraphics[width=\textwidth]{error_by_group}\\[2pt]
  \includegraphics[width=0.46\textwidth]{error_by_price_decile}\hfill
  \includegraphics[width=0.46\textwidth]{error_tier_distance}
  \caption{Error analysis on the test split. Top: Task~A MAE of
  $\log(\text{price})$ by category, by number of training listings in the
  district, and by true-price decile (row counts in brackets). Bottom left:
  mean residual per price decile; the tree pulls the extremes towards the
  mean. Bottom right: Task~B error rate vs.\ distance from the threshold.
  A map of the residuals is in the appendix (Fig.~\ref{fig:resmap}).}
  \label{fig:err}
\end{figure*}

\begin{figure}[t]
  \centering
  \includegraphics[width=\linewidth]{svm_weights}
  \caption{Largest weights of the final linear SVM (standardised
  features; positive values push towards premium).}
  \label{fig:weights}
\end{figure}

\begin{table}[t]
  \centering
  \caption{Largest Task~A test errors of our tree.}
  \label{tab:worst}
  \fit{\gen{tab_worst}}
\end{table}

\subsection{Bonus: ridge, ensembles of our trees, clustering}
\emph{Ridge} (normal equations, $\alpha=\RidgeAlpha$ chosen on
validation) matches scikit-learn's \texttt{Ridge} to machine precision.
\emph{Ensembles} (Fig.~\ref{fig:ens}, Table~\ref{tab:bonus}): a random
forest of \RFTrees{} of our trees (bootstrap + a third of the features per
split) \ForestVsTreeWord{} the test RMSE: \RFRMSE{} against \OurTreeRMSE{}
for a single tree (scikit-learn's forest: \SkRFRMSE). Least-squares gradient boosting with
\GBTrees{} depth-4 trees reaches \GBRMSE{} (scikit-learn: \SkGBRMSE).
The boosting curve \BoostCurveVerdict{} as learners are added. For Task~B
our forest reaches F1 \RFClfFOne{} (scikit-learn \SkRFClfFOne).
\emph{Clustering:} k-means++ ($k=\KMeansK$) on the projected coordinates
recovers neighbourhoods (Fig.~\ref{fig:km}). Per-cluster test MAE ranges
from \BestClusterMAE{} to \WorstClusterMAE{} (cluster \WorstClusterID).
Adding cluster indicators to the tree \AblationWord{} the validation
RMSE (\AblationWithout{} $\to$ \AblationWith). PCA needs \PCANinety{}
components for 90\% of the variance; the first two hold only
\PCATwoShare.

\begin{table}[t]
  \centering
  \caption{Bonus models on Task A (test split).}
  \label{tab:bonus}
  \fit{\gen{tab_bonus}}
\end{table}



% =====================================================================
\section{Discussion and Limitations}
\emph{Bias and variance.} A single deep tree has low bias and high
variance. Validation-chosen depth and leaf size cut the variance, and
averaging many decorrelated trees (the forest) cuts it further,
\EnsembleVerdict. The linear SVM has the opposite
profile: it cannot express interactions such as ``large \emph{and}
central'', and RFF buys that flexibility at the price of tuning $\gamma$.

\emph{Honest limitations.}
(1) Asking prices are not transaction prices, and re-posted listings
with edited prices can survive de-duplication.
(2) Cleaning thresholds and the EUR rate are judgement calls. They are
fixed and documented but not tuned, and filtering on price per m$^2$
removes some genuine extremes along with typos.
(3) One snapshot, one random split: CV quantifies sampling noise but not
temporal drift. A time-based split would be the stricter test.
(4) Task~B is a median split of Task~A's target, so it inherits all of
Task~A's noise and has an irreducible error band at the threshold.
(5) Our tree is several times slower than scikit-learn's, and Pegasos
with very small $\lambda$ needs many epochs. The RFF model is a
Monte-Carlo approximation whose result depends on the random seed.
(6) The description keywords are crude. A proper text model could add
signal, but it would also need care to avoid prices written in the text.

\section{Conclusion}
Both from-scratch models behave as specified: the tree matches
scikit-learn's split at every node up to ties and float32 effects, and
Pegasos reaches LIBLINEAR's objective up to \ObjGap. Both are
competitive on real, messy data. Most of the work, and most of the error, lies in the data: currency
quirks, duplicates, sparse districts and the luxury tail. Next steps are
a time-based split, district-level target encoding learned inside the
folds, quantile losses for the tails, and histogram-based split finding
to close the speed gap.

\section*{Contributions}
See \texttt{contribution\_report.md} in the repository root for who did
what.

\section*{Tools and Acknowledgements}
Data: the Kaggle \emph{Bina.az sale project} dataset~\cite{data}; we do
not redistribute it. scikit-learn~\cite{sklearn} is used only for the
baselines, the bonus comparisons and the unit tests that check our
metrics and models. \textbf{AI assistance:} an AI coding assistant
(Claude Code, Anthropic) was used to draft and refactor parts of the
code, tests and this report. All of it was reviewed, run and checked by
the team, and every number in this report is regenerated by
\texttt{python -m src.run\_all}. % TEAM: make this statement accurate for your own use.

% =====================================================================
\begin{thebibliography}{99}\small
\bibitem{pegasos}
S.~Shalev-Shwartz, Y.~Singer, N.~Srebro, and A.~Cotter,
``Pegasos: primal estimated sub-gradient solver for SVM,''
\emph{Math.\ Programming}, vol.~127, no.~1, pp.~3--30, 2011.
\bibitem{cart}
L.~Breiman, J.~Friedman, R.~Olshen, and C.~Stone,
\emph{Classification and Regression Trees}. Wadsworth, 1984.
\bibitem{quinlan}
J.~R.~Quinlan, ``Induction of decision trees,''
\emph{Machine Learning}, vol.~1, no.~1, pp.~81--106, 1986.
\bibitem{cortes}
C.~Cortes and V.~Vapnik, ``Support-vector networks,''
\emph{Machine Learning}, vol.~20, no.~3, pp.~273--297, 1995.
\bibitem{rahimi}
A.~Rahimi and B.~Recht, ``Random features for large-scale kernel
machines,'' in \emph{Proc.\ NeurIPS}, 2007, pp.~1177--1184.
\bibitem{lacoste}
S.~Lacoste-Julien, M.~Schmidt, and F.~Bach, ``A simpler approach to
obtaining an $O(1/t)$ convergence rate for the projected stochastic
subgradient method,'' arXiv:1212.2002, 2012.
\bibitem{bottou}
L.~Bottou, ``Large-scale machine learning with stochastic gradient
descent,'' in \emph{Proc.\ COMPSTAT}, 2010, pp.~177--186.
\bibitem{liblinear}
R.-E.~Fan, K.-W.~Chang, C.-J.~Hsieh, X.-R.~Wang, and C.-J.~Lin,
``LIBLINEAR: A library for large linear classification,''
\emph{JMLR}, vol.~9, pp.~1871--1874, 2008.
\bibitem{sklearn}
F.~Pedregosa \emph{et al.}, ``Scikit-learn: Machine learning in Python,''
\emph{JMLR}, vol.~12, pp.~2825--2830, 2011.
\bibitem{breiman}
L.~Breiman, ``Random forests,'' \emph{Machine Learning}, vol.~45, no.~1,
pp.~5--32, 2001.
\bibitem{friedman}
J.~H.~Friedman, ``Greedy function approximation: a gradient boosting
machine,'' \emph{Ann.\ Statist.}, vol.~29, no.~5, pp.~1189--1232, 2001.
\bibitem{kmeanspp}
D.~Arthur and S.~Vassilvitskii, ``k-means++: the advantages of careful
seeding,'' in \emph{Proc.\ ACM-SIAM SODA}, 2007, pp.~1027--1035.
\bibitem{data}
sehriyarmemmedli, ``Bina.az sale project,'' Kaggle dataset.
\url{https://www.kaggle.com/datasets/sehriyarmemmedli/binaaz-sale-project}
\end{thebibliography}

% =====================================================================
\clearpage
\onecolumn
\appendix
\section*{Appendix: additional figures}

\begin{figure}[h]
  \centering
  \includegraphics[width=0.48\textwidth]{eda_location_effect}\hfill
  \includegraphics[width=0.48\textwidth]{eda_premium_share}\\[4pt]
  \includegraphics[width=0.48\textwidth]{eda_price_by_rooms}
  \caption{Location effect (median price per m$^2$ by district, training
  listings in brackets), premium share by category, and $\log$ price by
  room count (training split).}
  \label{fig:loc}
\end{figure}

\begin{figure}[h]
  \centering
  \includegraphics[width=0.48\textwidth]{error_residual_map}
  \caption{Test residuals of our Task~A tree across Baku (red: over-predicted,
  blue: under-predicted).}
  \label{fig:resmap}
\end{figure}

\begin{figure}[h]
  \centering
  \includegraphics[width=\textwidth]{tree_param_sweeps}\\[4pt]
  \includegraphics[width=\textwidth]{tree_learning_curves}
  \caption{Top: validation RMSE vs.\ \texttt{min\_samples\_leaf} for a
  fully grown tree, and vs.\ \texttt{min\_impurity\_decrease} at the
  selected depth. Bottom: learning curves of the selected trees.}
  \label{fig:tree-extra}
\end{figure}

\begin{figure}[h]
  \centering
  \includegraphics[width=\textwidth]{results_roc_pr}\\[4pt]
  \includegraphics[width=\textwidth]{results_confusion}\\[4pt]
  \includegraphics[width=\textwidth]{results_pred_vs_actual}
  \caption{Test-set ROC and precision--recall curves, confusion matrices,
  and predicted vs.\ actual $\log$ price.}
  \label{fig:curves}
\end{figure}

\begin{figure}[h]
  \centering
  \includegraphics[width=\textwidth]{bonus_ensembles}
  \caption{Bonus: test RMSE as learners are added to our random forest (a) and our gradient
  boosting (b), with scikit-learn's final value and our single tree as references.}
  \label{fig:ens}
\end{figure}

\begin{figure}[h]
  \centering
  \includegraphics[width=\textwidth]{bonus_kmeans}\\[4pt]
  \includegraphics[width=\textwidth]{bonus_pca}
  \caption{Bonus: k-means elbow and geographic clusters coloured by
  median $\log$ price; PCA explained variance and the two tiers in the
  first two components.}
  \label{fig:km}
\end{figure}

\end{document}
